\section{Definición del problema}\label{sec:problema}

\subsection{Pregunta de investigación}
\begin{quote}
\itshape ¿Cómo cambiaron los patrones de movilidad, seguridad y actividad económica de la
Ciudad de México durante la pandemia de COVID-19, y qué diferencias se observan en su
recuperación posterior respecto al periodo previo a 2020?
\end{quote}

\subsection{Planteamiento}
La pandemia de COVID-19 provocó un choque abrupto sobre la vida urbana de la Ciudad de México (CDMX).
En México, el primer caso se confirmó el 28 de febrero de 2020 y la Jornada Nacional de Sana Distancia
comenzó el 23 de marzo de 2020.

Los datos para medir este choque existen, pero están \textbf{dispersos} entre distintas instituciones:
Gobierno de la CDMX, Fiscalía General de Justicia (FGJ), Secretaría de Seguridad Ciudadana (SSC), C5,
INEGI, IMSS y Secretaría de Salud. Además, tienen \textbf{granularidades distintas}. En el tiempo
pueden ser diarios, mensuales o trimestrales. En el territorio pueden estar por estación, colonia,
alcaldía o entidad. Sus formatos y calidades también varían. No existe una base integrada que permita
comparar las tres dimensiones sobre un mismo eje temporal y territorial.

\begin{table}[h]
\centering
\caption{Problemas de los datos disponibles.}
\small
\begin{tabular}{L{3cm}L{11cm}}
\toprule
\textbf{Problema} & \textbf{Ejemplo}\\
\midrule
Dispersión & Gobierno CDMX, FGJ, SSC, C5, INEGI, IMSS y Secretaría de Salud publican cada uno por separado.\\
Granularidad distinta & Metro: diario por estación. Carpetas de la FGJ: por evento, con coordenadas. ITAEE: trimestral por entidad.\\
Calidad heterogénea & Texto mal codificado, cambios de catálogo (PGJ $\rightarrow$ FGJ), cierres de líneas del Metro y series que empiezan en 2020.\\
\bottomrule
\end{tabular}
\end{table}

\subsection{Qué se analizará}
El análisis cubre tres dimensiones de la vida urbana, más una dimensión de control que mide la
intensidad de la pandemia (Tabla~\ref{tab:dimensiones}).

\begin{table}[h]
\centering
\caption{Dimensiones y variables del análisis.}\label{tab:dimensiones}
\small
\begin{tabular}{L{2.6cm}L{3.8cm}L{7.6cm}}
\toprule
\textbf{Dimensión} & \textbf{Qué se mide} & \textbf{Variables principales}\\
\midrule
Movilidad & Uso del transporte público y privado & Afluencia de Metro, Metrobús, STE, RTP y Ecobici; tránsito vehicular; hechos de tránsito\\
Seguridad & Delitos denunciados, emergencias y percepción & Carpetas de investigación por tipo de delito, víctimas, llamadas al 911, percepción de inseguridad (ENSU)\\
Actividad económica & Empleo y actividad productiva & Asegurados en el IMSS, ocupación y desocupación (ENOE), ITAEE, unidades económicas (DENUE), ocupación hotelera\\
Pandemia (control) & Intensidad del COVID-19 & Casos, positividad, hospitalización, color del semáforo\\
\bottomrule
\end{tabular}
\end{table}

\subsection{Objetivo de la fase}
El problema a resolver en esta fase es \textbf{construir un conjunto de datos integrado, limpio y
documentado} que permita:
\begin{enumerate}
  \item Medir la magnitud de la caída de cada dimensión en 2020.
  \item Medir la velocidad y el nivel de recuperación: si se volvió a la línea base, cuándo, o si se
        llegó a un nuevo equilibrio.
  \item Relacionar las dimensiones entre sí y con la intensidad de la pandemia (casos, semáforo epidemiológico).
  \item Servir de insumo a un modelo predictivo en la fase siguiente.
\end{enumerate}

\subsection{Periodos de análisis}
\begin{table}[h]
\centering
\caption{Periodos propuestos. Los cortes pueden ajustarse con base en el semáforo epidemiológico de la CDMX.}
\begin{tabular}{lll}
\toprule
\textbf{Periodo} & \textbf{Rango} & \textbf{Uso} \\
\midrule
Pre-pandemia (línea base) & 2016-01-01 -- 2020-02-29 & Tendencia y estacionalidad \enquote{normal} \\
Confinamiento / choque    & 2020-03-01 -- 2021-06-30 & Caída; semáforo rojo/naranja \\
Transición                & 2021-07-01 -- 2022-12-31 & Reapertura gradual, vacunación \\
Post-pandemia             & 2023-01-01 -- último dato & Recuperación / \enquote{nueva normalidad} \\
\bottomrule
\end{tabular}
\end{table}

\subsection{Unidades de análisis}
\begin{itemize}
  \item \textbf{Temporal:} diaria donde exista. La base integrada se agregará a nivel \textbf{semanal y mensual}.
  \item \textbf{Territorial:} la unidad común es la \textbf{alcaldía} (16). Se usará colonia o estación
        solo cuando el dato lo permita.
\end{itemize}

\subsection{Variables objetivo candidatas (fase predictiva)}
\begin{itemize}
  \item Afluencia mensual total en transporte público (Metro + Metrobús + STE + RTP).
  \item Delitos de alto impacto por alcaldía y mes.
  \item Empleo formal (trabajadores asegurados en el IMSS) en la CDMX por mes.
  \item Índice de recuperación: el valor observado entre el valor esperado sin pandemia. El valor
        esperado es un contrafactual estimado con la tendencia anterior a 2020.
\end{itemize}

\subsection{Hipótesis de trabajo}
\begin{description}
  \item[H1.] La movilidad en transporte masivo cayó más del 50\,\% en abril--mayo de 2020 y no se
             recuperó al mismo ritmo que el tránsito vehicular.
  \item[H2.] Los delitos patrimoniales en vía pública (robo a transeúnte y en transporte) cayeron junto
             con la movilidad. La violencia familiar no disminuyó, o incluso aumentó.
  \item[H3.] El empleo formal se recuperó antes que la movilidad en transporte público, por el efecto
             del teletrabajo.
  \item[H4.] La recuperación fue desigual entre alcaldías centrales y periféricas. Con la selección
             final (Sección~\ref{sec:seleccion_final}) solo se evalúa en seguridad y seguridad vial.
\end{description}

\subsection{Alcance y limitaciones}
\begin{itemize}
  \item El modelado, la inferencia causal y los pronósticos quedan fuera de esta fase.
  \item Las carpetas de investigación miden \emph{denuncias}, no delitos ocurridos. Según la ENVIPE,
        la cifra negra ronda el 90\,\%. Por eso se complementan con las llamadas al 911 y con las encuestas ENVIPE y ENSU.
  \item Algunas series creadas durante la pandemia (Google Mobility, índices de movilidad de la CDMX)
        \textbf{no tienen línea base anterior a 2020}. Sirven para describir el periodo pandémico, pero
        no para comparar antes y después.
  \item Hubo cambios metodológicos. Por ejemplo, la PGJ se transformó en FGJ y algunos delitos se
        reclasificaron. Estos cambios deben documentarse al unir las series.
\end{itemize}
